% Einheit 2 von 5 – Prozentsatz berechnen (bank/prozentrechnung/e2.jsonl, zusammenbau v0.1)
%% TODO Zweigzeile Teil 1: was hier gelernt wird (Satz in Schülersprache aus dem Einheitstitel) – Einheit 2
\einheitenkopf[e2]{Einheit 2 von 5 · Prozentsatz berechnen}
\zweigzeile{neu in diesem Jahr · P10 oft}
\setcounter{aufgabe}{11}

%% TODO Titel: Ich-kann-Satz fehlt in der Bank (Platzhalter: Kettenname) – Nr. 12
%% TODO Anweisung: ein Satz über der ersten Teilaufgabe fehlt in der Bank – Nr. 12
\begin{aufgabe}{Was ist das Ganze?}
\begin{teile}
\teil In der Klasse sind $28$ Kinder, $7$ davon haben einen Hund. – Was ist das Ganze?
\end{teile}
\end{aufgabe}

%% TODO Titel: Ich-kann-Satz fehlt in der Bank (Platzhalter: Kettenname) – Nr. 13
%% TODO Anweisung: ein Satz über der ersten Teilaufgabe fehlt in der Bank – Nr. 13
\begin{aufgabe}{Streifen einteilen}
\swa{In $10$-\%-Schritte einteilen.}{\streifenleer[0]}
\end{aufgabe}

%% TODO Titel: Ich-kann-Satz fehlt in der Bank (Platzhalter: Kettenname) – Nr. 14
%% TODO Anweisung: ein Satz über der ersten Teilaufgabe fehlt in der Bank – Nr. 14
\begin{aufgabe}{Prozentsatz}
\swa{Streifen: $50$ Kinder, grau: $10$ Kinder – wie viel Prozent? \leerfeld[\%]}{\streifen{20}{$0$}{$50$ Kinder}}
%% TODO Darstellung für schwach fehlt in der Bank (prozentrechnung-e2-k3-s1-v1; Abschnitt „Für schwache Schüler“)
\swz{$17$ von $100$ Schülern spielen ein Instrument – wie viel Prozent?}{}
%% TODO Darstellung für schwach fehlt in der Bank (prozentrechnung-e2-k3-s1-v2; Abschnitt „Für schwache Schüler“)
\swz{$68$ von $100$ Losen sind Nieten – wie viel Prozent?}{}
%% TODO Darstellung für schwach fehlt in der Bank (prozentrechnung-e2-k3-s1-v3; Abschnitt „Für schwache Schüler“)
\swz{$3$ von $100$ Äpfeln sind faul – wie viel Prozent?}{}
%% TODO Darstellung für schwach fehlt in der Bank (prozentrechnung-e2-k3-s1-v4; Abschnitt „Für schwache Schüler“)
\swz{$45$ von $100$ Befragten fahren Rad – wie viel Prozent?}{}
%% TODO Darstellung für schwach fehlt in der Bank (prozentrechnung-e2-k3-s1-v5; Abschnitt „Für schwache Schüler“)
\swz{$100$ Gäste, $81$ davon bleiben zum Essen – wie viel Prozent?}{}
\swfrage{Was bleibt gleich, was ändert sich?}
%% TODO Darstellung für schwach fehlt in der Bank (prozentrechnung-e2-k3-s2-v1; Abschnitt „Für schwache Schüler“)
\swz{$19$ von $50$ Kindern – wie viel Prozent?}{}
%% TODO Darstellung für schwach fehlt in der Bank (prozentrechnung-e2-k3-s3-v1; Abschnitt „Für schwache Schüler“)
\swz{$27$ von $60$ Minuten Training sind Laufen – wie viel Prozent?}{}
%% TODO Darstellung für schwach fehlt in der Bank (prozentrechnung-e2-k3-s4-v1; Abschnitt „Für schwache Schüler“)
\swz{$17$ von $24$ Kindern waren im Schwimmbad – wie viel Prozent? Runde auf eine Stelle.}{}
%% TODO Darstellung für schwach fehlt in der Bank (prozentrechnung-e2-k3-s5-v1; Abschnitt „Für schwache Schüler“)
\swz{Im Chor singen $14$ Mädchen und $6$ Jungen – wie viel Prozent sind Jungen?}{}
%% TODO Darstellung für schwach fehlt in der Bank (prozentrechnung-e2-k3-s6-v1; Abschnitt „Für schwache Schüler“)
\swz{Jacke: vorher $80\,€$, jetzt $60\,€$ – wie viel Prozent Rabatt?}{}
\swz[3]{Bei einer Tombola gibt es $21$ Nieten und $3$ Hauptgewinne. Wie viel Prozent der Lose sind Hauptgewinne? \hfill (P10 2018 OS)}{}
\swz[3]{Eine Stadt hat $40\,000$ Einwohner, $1\,080$ sind in diesem Jahr zugezogen. Wie viel Prozent sind das? \hfill (P10 2014 OS)}{}
\swz[3]{Körpergrößen der $13$ Spieler einer Mannschaft in cm: $172$, $185$, $169$, $190$, $178$, $181$, $166$, $188$, $175$, $183$, $170$, $192$, $177$. Wie viel Prozent sind größer als $180$ cm? Runde auf eine Stelle. \hfill (P10 2025 OS)}{}
\end{aufgabe}

%% TODO Titel: Ich-kann-Satz fehlt in der Bank (Platzhalter: Kettenname) – Nr. 15
%% TODO Anweisung: ein Satz über der ersten Teilaufgabe fehlt in der Bank – Nr. 15
\begin{aufgabe}{Umkehrung: zu einem Prozentsatz ein Zahlenpaar angeben}
%% TODO Darstellung für schwach fehlt in der Bank (prozentrechnung-e2-k4-s1-v1; Abschnitt „Für schwache Schüler“)
\swz{$30\,\%$ – Teil und Ganzes?}{}
\end{aufgabe}

%% TODO Titel: Ich-kann-Satz fehlt in der Bank (Platzhalter: Kettenname) – Nr. 16
%% TODO Anweisung: ein Satz über der ersten Teilaufgabe fehlt in der Bank – Nr. 16
\begin{aufgabe}{Prozentsatz – Fehler finden, Begründen, Darstellungswechsel, Anwendung}
\swz[3]{Kim, „$12$ von $48$ Kindern“: \rechnung{48 : 12 &= 4 = 400\,\%} Fehler? Wie heißt es richtig?}{}
\swfrage{Warum teilt man beim Prozentsatz durch das Ganze und nicht durch den Teil?}
\swa{Ganzes $40$ kg, Teil $10$ kg – im Streifen einzeichnen, Prozentsatz ablesen. \leerfeld[\%]}{\streifenleer[4]}
\swz[3]{Handyvertrag mit $20$ GB im Monat. Am 20. Tag sind $17$ GB verbraucht. Wie viel Prozent sind noch übrig? Reicht der Rest bei gleichem Verbrauch für die letzten $10$ Tage?}{}
\end{aufgabe}

\uebersichtskasten{Prozentsatz: Teil geteilt durch Ganzes – das ist der Anteil, dann in Prozent schreiben. \\ 14 von 40: 14 : 40 = 0,35 = 35 \% \quad 3 von 8: 3 : 8 = 0,375 = 37,5 \%}
